% Generated from locale/tr/content/computability/computability-theory/computable-sets.tex; do not hand-edit.


\olfileid[tr]{cmp}{thy}{cps}
\olsection{Hesaplanabilir Kümeler}

Hesaplanabilirlik kavramını hesaplanabilir fonksiyonlardan hesaplanabilir
kümelere genişletebiliriz.

\begin{defn}
  $S$ bir doğal sayılar kümesi olsun. $S$, ancak ve ancak karakteristik
  fonksiyonu $\Char{S}$ hesaplanabilirse \emph{hesaplanabilir}dir. Başka bir
  deyişle $S$, ancak ve ancak
  \[
    \Char{S}(x)=\begin{cases}
      1 & \text{$x\in S$ ise}\\
      0 & \text{aksi durumda}
    \end{cases}
  \]
  fonksiyonu hesaplanabilirse hesaplanabilirdir. Benzer biçimde bir
  $R(x_0,\dots,x_{k-1})$ bağıntısı ancak ve ancak karakteristik fonksiyonu
  hesaplanabilirse hesaplanabilirdir.

  Hesaplanabilir küme ve bağıntılara \emph{karar verilebilir} de denir.
\end{defn}

\begin{explain}
Artık kısmi fonksiyonlar, fonksiyonlar ve kümeler için ayrı hesaplanabilirlik
kavramlarımız bulunduğuna dikkat edin; bunları karıştırmayın.
\iftag{TMs}{Kısmi bir fonksiyonu hesaplayan Turing makinesi, fonksiyonun
tanımlı olduğu girdi değerlerinde çıktıyı döndürür. Bir kümeyi hesaplayan
Turing makinesi ise girdi değerinin kümede bulunup bulunmadığına karar
verdikten sonra $1$ ya da $0$ döndürür.}{}
\end{explain}

